\par\Needspace{9\baselineskip}\bigskip
\pdfbookmark[3]{Exposition, [25]}{pallas-II-article25}
\begin{center}[25.]\end{center}
\begingroup
\noindent\hspace*{\parindent}Nous indiquerons à présent une méthode très simple, de calculer complètement les perturbations de l’ordre du quarré de la force perturbatrice. Soit \(Q\mathrm{d}t\) la variation instantanée d’un élément, p. e. de \(\AscendingNode+\delta\AscendingNode\); \(Q\) dépendra des élémens de la planète perturbée et de la planète perturbante. Si l’on y substitue les valeurs de ces élémens justes au premier ordre des perturbations, il est clair que \(Q\) sera juste au second ordre (inclusiv.). Soient donc \SourceLeaf{475}\(a,\ b,\ c,\ e,\ f,\ g\) les valeurs moyennes des élémens de la planète perturbée, \(a+\Delta a,\ b+\Delta b,\ c+\Delta c,\ e+\Delta e,\ f+\Delta f,\ g+\Delta g\) les valeurs justes au premier ordre; et que \(a',\ b',\ c',\ e',\ f',\ g',\ a'+\Delta a'\) etc. désignent les mêmes quantités pour la planète perturbante. On pourra supposer
\[
\begin{alignedat}{2}
Q=Q^0&{}+\Delta a\cdot\left(\frac{\mathrm{d}Q^0}{\mathrm{d}a}\right)&&{}+\Delta a'\cdot\left(\frac{\mathrm{d}Q^0}{\mathrm{d}a'}\right)\\[3pt]
&{}+\Delta b\cdot\left(\frac{\mathrm{d}Q^0}{\mathrm{d}b}\right)&&{}+\Delta b'\cdot\left(\frac{\mathrm{d}Q^0}{\mathrm{d}b'}\right)\\[3pt]
&{}+\Delta c\cdot\left(\frac{\mathrm{d}Q^0}{\mathrm{d}c}\right)&&{}+\Delta c'\cdot\left(\frac{\mathrm{d}Q^0}{\mathrm{d}c'}\right)\\[3pt]
&\quad\text{etc.}&&\quad\text{etc.,}
\end{alignedat}
\]
\noindent si l’on se borne aux quantités de second ordre. Les quantités \(\Delta a,\allowbreak\ \Delta b,\allowbreak\ \Delta c,\allowbreak\ \Delta e,\allowbreak\ \Delta f,\allowbreak\ \Delta g\) contiennent chacune un terme proportionnel au tems et des fonctions périodiques de \(nt\) et \(n't\) ou .\ .\ .\ .\ .
\par\medskip\begin{center}\rule{.18\linewidth}{.3pt}\end{center}
\endgroup
